\clearpage
\section{Chapter 4: Standard Model}

\subsection{Questions}

\begin{enumerate}[label=4.\arabic*]
  \item \label{q:4.1} \textbf{Particle Classification and Properties}

  For each particle below, identify whether it is a fermion or boson, its typical spin, and which interaction(s) it directly mediates and/or participates in:
  \[
    e^-,\quad \nu_e,\quad \gamma,\quad g,\quad H.
  \]
  (Here $g$ denotes gluon and $H$ denotes Higgs boson.)

  \item \label{q:4.2} \textbf{Hadrons and Quark Content}

  \begin{enumerate}[label=(\alph*)]
    \item Classify each as baryon or meson: $p$, $n$, $\pi^+$, $K^0$.
    \item Write quark content for $p$, $n$, and $\pi^+$.
    \item Why are free quarks not observed at low energy scales?
  \end{enumerate}

  \item \label{q:4.3} \textbf{Higgs Mechanism and Degrees of Freedom}

  Use the notes relation
  \[
    N_{\mathrm{dof}} = 2N_G + N_H = 3N_G + (N_H-N_G),
  \]
  where $N_G$ is number of gauge bosons that become massive, and $N_H$ is number of Higgs-field real DoF.
  For electroweak symmetry breaking, take $N_G=3$ ($W^\pm,Z$) and $N_H=4$.
  \begin{enumerate}[label=(\alph*)]
    \item Compute total DoF before symmetry breaking.
    \item Compute total DoF after symmetry breaking.
    \item Interpret physically what the leftover $(N_H-N_G)$ DoF corresponds to.
  \end{enumerate}
\end{enumerate}

\bigskip
\bigskip
\bigskip

\subsection{Solutions}

\subsubsection*{Solution 4.1: Particle Classification and Properties}

\begin{itemize}
  \item $e^-$: fermion, spin $1/2$, charged lepton (participates in EM and weak interactions).
  \item $\nu_e$: fermion, spin $1/2$, neutrino (weak interaction; neutral under EM and color).
  \item $\gamma$: boson, spin $1$, gauge boson of electromagnetism.
  \item $g$ (gluon): boson, spin $1$, gauge boson of strong interaction; carries color.
  \item $H$ (Higgs): boson, spin $0$, scalar responsible for electroweak symmetry breaking and mass generation via Higgs mechanism.
\end{itemize}

\subsubsection*{Solution 4.2: Hadrons and Quark Content}

\begin{enumerate}[label=(\alph*)]
  \item
  \[
    p, n:\ \text{baryons};\qquad \pi^+, K^0:\ \text{mesons}.
  \]
  (Baryon = 3 quarks, meson = quark + antiquark.)

  \item
  \[
    p = uud,\qquad n = udd,\qquad \pi^+ = u\bar d.
  \]

  \item At low energies, QCD confinement means strong coupling becomes large; isolating quarks costs increasing energy, so free quarks are not observed.
\end{enumerate}

\subsubsection*{Solution 4.3: Higgs Mechanism and Degrees of Freedom}

\begin{enumerate}[label=(\alph*)]
  \item Before symmetry breaking:
  \[
    N_{\mathrm{dof}} = 2N_G + N_H = 2\times 3 + 4 = 10.
  \]

  \item After symmetry breaking:
  \[
    N_{\mathrm{dof}} = 3N_G + (N_H-N_G)=3\times3 + (4-3)=10.
  \]
  Total DoF is conserved.

  \item The $N_H-N_G = 1$ leftover scalar DoF corresponds to the physical Higgs boson.
\end{enumerate}


% ─────────────────────────────────────────────────────────────────────────────
